\begin{python0}
from solutions import *; clear()
\end{python0}
\sectionthree{Read Error}

Now suppose you've checked that the file is open for reading, what can happen?

You might be attempting reading an integer from file \texttt{f}:

\begin{console}
int x;
f >> x;
\end{console}

when the contents from \texttt{f} to be processed is not an integer. For instance from the file pointer, the next sequence of data is \texttt{"hello"}. This is sometimes called a format error.

Such an error can be detected like this:
\begin{console}
int x;
f >> x;
if (f.fail())
{  
    std::cout <<"x is garbage\n";
    char y[100];
    f >> y;
    std::cout << "data read: " << y << '\n';
}
else
{   
    // do something useful with x
}
\end{console}

In this case data is not lost. In other words, the file pointer does not
move:

\begin{consolethree}[escapeinside=||]
int x;
f >> x;
if (f.fail())
{    
     std::cout << "x is garbage\n";
     |\textbf{char y[100];}|
     |\textbf{f >> y;}|
     |\textbf{std::cout << " data read: " << y << '\textbackslash n';}|
}
else
{    
     // do something useful with x
}
\end{consolethree}

There are other situations where error is more serious (example:
hardware failure) where data can be lost. This is how you check for such
errors:

\begin{console}
int x;
f >> x;
if (f.bad())
{  
    std::cout << "fatal error\n";
}
else
{   
    // do something useful with x
}
\end{console}

Such cases are really bad and recovery is difficult.

A detailed error handling code (assuming the file is already checked to be open) might look like this:

\begin{console}
f >> x;
if (f.fail())
{   
    // try to do some recovery, maybe read x in a
    // different way and proceed
}
else if (f.bad())
{   
    // fatal error ... maybe print error message and
    // halt the program
}
else
{   
    // everything is OK with reading x
}
\end{console}

If you don't care to handle \texttt{f.fail()} and
\texttt{f.bad()} cases separately you can of course do this:

\begin{console}
f >> x;
if (f.fail() || f.bad())
{   
    // don't bother with recovery ... data was bad
    // anyway ... GIGO
    // just print an error message and halt the
    // program
}
else
{   
    // everything is OK with reading x ... MOVE ON!!!
}
\end{console}

There's a shorthand for doing the above.

The opposite of \texttt{f.bad()} is \texttt{f.good()}. So you can do this if you like:

\begin{console}
f >> x;
if (f.good())
{    
     // everything is OK with reading x ... MOVE ON!!!
}
else
{   
     // don't bother with recovery ... data was bad
     // anyway ... GIGO
     // just print an error message and halt the
     // program
} 
\end{console}

So let's say you have a file containing first name, last name, monthly pay:

\begin{console}
John Doe 1234.56
Jane Smith 2345.67
Harry Lee 3456.78
\end{console}

You can do this (assuming you have already checked that your file \texttt{f} is opened successfully):

\begin{console}
double total = 0.0;
bool error = false;
while (!f.eof())
{   
     char[100] fname;
     char[100] lname;
     double pay;
     f >> fname >> lname >> pay;
     if (f.good())
     {
          total += pay;
          std::cout << fname << ' ' << lname
                    << " is paid "
                    << pay << std::endl;
}
else
{    
     error = true;
     break;
}
if (!error) // or use f.good()
{    
     std::cout << "total: " << total << std::endl;
}
else
{
std::cout << "DATA IN FILE IS CORRUPTED!!!"
          << std::endl;
}
\end{console}

